% Generated draft: requires mathematical and source review.
\documentclass[11pt]{article}
\usepackage[margin=25mm]{geometry}
\usepackage{fontspec}
\setmainfont{Latin Modern Roman}
\usepackage{amsmath,amssymb,mathtools,unicode-math}
\setmathfont{Latin Modern Math}
\usepackage{enumitem,needspace,xurl,hyperref}
\setlength{\parindent}{0pt}
\setlength{\parskip}{5pt}
\setlength{\emergencystretch}{3em}
\newcommand{\N}{\mathbb{N}}
\newcommand{\Z}{\mathbb{Z}}
\newcommand{\Q}{\mathbb{Q}}
\newcommand{\R}{\mathbb{R}}
\newcommand{\C}{\mathbb{C}}
\newcommand{\F}{\mathbb{F}}
\newcommand{\A}{\mathbb{A}}
\newcommand{\PP}{\mathbb{P}}
\newcommand{\E}{\mathbb{E}}
\newcommand{\eps}{\varepsilon}
\newcommand{\dd}{\,\mathrm{d}}
\begin{document}
\section*{UnsolvedMath ID: 20000445}
\textbf{Source problem code:} catalogue entry
\subsection{Definitions and mathematical statement}
The original Rodriguez-Villegas question asks to classify Galois-stable subsets \(A,B\subset\mu_\infty\) whose elements alternate around the unit circle, but its extracted condition says \(A=B\). The attached report repairs this explicitly to \(|A|=|B|\). For a precise strict-interlacing formulation take nonempty finite disjoint sets of equal size, stable under \(\operatorname{Gal}(\overline\Q/\Q)\), whose cyclicly ordered union alternates between the two sets. Literal equality of the sets is incompatible with that strict interpretation and is not silently retained as an equivalent condition.

Galois stability means that each set is a union of primitive-root orbits, equivalently its root polynomial is a product of distinct cyclotomic polynomials. The report attributes the known classification to Beukers--Heckman in the explicit McKee--Smyth form: two primitive infinite families and \(26\) primitive sporadic pairs, with all pairs obtained up to interchange, simultaneous sign change, and common inflation. The second family, sporadic list, and precise primitiveness conventions are not supplied, so this record does not contain a complete reproducible classification.

The original remark gives one infinite family, for positive coprime integers \(m,n\),
\[
A=\mu_{m+n}\setminus\{1\},\qquad B=\mu_m\cup\mu_n.
\]
It says a solution is known but more complicated than the statement and is used in classifying algebraic hypergeometric functions.

The report's candidate discrepancy certificate can be stated after orienting \(1\in B\). Set \(\epsilon_d=1\) for primitive \(d\)-th-root orbits in \(A\), \(-1\) for those in \(B\), and zero otherwise. With \(\mu(e)\) the number-theoretic M\"obius function, define, for \(0\leq x<1\),
\[
\Delta_\epsilon(x)=\sum_d\epsilon_d\sum_{e\mid d}\mu(e)\left\lfloor\frac{dx}{e}\right\rfloor.
\]
The sum has finite support. In these conventions it counts the signed nonzero-angle roots encountered from angle zero through \(2\pi x\), with \(1\in B\) as the initial endpoint. Under the finite, disjoint, equal-cardinality orbit hypotheses, the report asserts strict interlacing is equivalent to
\[
\Delta_\epsilon(x)\in\{0,1\}\quad\text{for every }0\leq x<1.
\]
The endpoint and sign convention matter for this exact range.

For common inflation by an integer \(\ell\geq1\), meaning \(A_\ell=\{z:z^\ell\in A\}\) and \(B_\ell=\{z:z^\ell\in B\}\), the candidate exact law is
\[
\Delta_{\mathrm{inflated}}(x)=\Delta_{\mathrm{base}}(\{\ell x\}),
\]
where braces mean fractional part. The report also claims one-bit floor-carry proofs of both classified infinite families, but does not give the second family's formula. The certificate and inflation identity are candidate additions; the underlying known classification is not claimed as newly established.
\subsection{Short English statement}
Classify finite Galois-stable root-of-unity sets that strictly interlace. The report repairs equality to equal cardinality, cites a known classification, and proposes an exact signed floor-sum certificate and inflation law while omitting the full classification list.
\subsection{Sources}
\begin{itemize}
\item AIM workshop ``Rational and integral points on higher dimensional varieties,'' Question 12, attributed to Fernando Rodriguez-Villegas: \url{https://aimath.org/WWN/qptsurface2/qptsurface2.pdf}. No PDF was accessed here.
\item Local canonical registry \texttt{aim-algebraic-number-theory-notes.json}, \texttt{notes[96]}, tag \texttt{question}; permanent UnsolvedMath ID \(20000445\), source code \texttt{catalogue entry}. The raw equality condition and the original infinite-family and hypergeometric-classification remarks are retained. Original literature field is empty.
\item Supplied research attempt 1 attributes the classification to Beukers--Heckman in McKee--Smyth's explicit form and proposes the discrepancy/inflation synthesis. No full authoritative references, second-family formula, or sporadic list are supplied. The report is machine generated, not peer reviewed, and requires independent verification before citation.
\end{itemize}
\end{document}
